\chapter{终端 UI：React + Ink}

\section{技术选型}

Claude Code 使用 \textbf{React + Ink} 构建终端 UI——这是一个非常规但极其优雅的选择。\href{https://github.com/vadimdemedes/ink}{Ink} 是一个使用 React 语法构建命令行界面的框架，将 React 的组件模型和声明式编程带到终端。

\section{为什么用 React 做终端 UI}

\begin{longtable}{ll}
\toprule
传统方案 & React + Ink \\
\midrule
命令式字符串拼接 & 声明式组件 \\
手动管理状态更新 & React 状态管理 \\
难以复用 UI 逻辑 & 组件化复用 \\
难以测试 & 可测试的组件 \\
难以实现复杂布局 & Flexbox 布局 \\
\bottomrule
\end{longtable}

\section{React Canary 版本}

\begin{lstlisting}[breaklines=true]
{
  "react": "^19.3.0-canary-705268dc-20260409",
  "react-reconciler": "^0.34.0-canary-705268dc-20260409"
}
\end{lstlisting}

Claude Code 使用 React 19 的 canary 版本。这是因为 Ink 需要与 \code{react-reconciler} 紧密配合，而 reconciler 的稳定版本可能不支持最新的 Ink 特性。

\section{组件库概览}

\code{src/components/} 目录包含约 \textbf{140 个组件}，覆盖了整个 CLI 界面：

\subsection{核心组件}

\begin{longtable}{ll}
\toprule
组件 & 功能 \\
\midrule
\code{REPL.tsx} & 主 REPL 界面——整个应用的顶层组件 \\
\code{MessageSelector.tsx} & 消息选择器（选择发送哪些消息给 API） \\
\code{AgentProgressLine.tsx} & 代理进度显示 \\
\code{Spinner.tsx} & 加载动画 \\
\bottomrule
\end{longtable}

\subsection{Ink 渲染器封装}

\begin{lstlisting}[language=TypeScript,breaklines=true]
src/ink/
├── renderer.ts        # 自定义 Ink 渲染器
└── termio/
    └── dec.ts         # 终端控制序列
\end{lstlisting}

\begin{lstlisting}[language=TypeScript,breaklines=true]
import { SHOW_CURSOR } from './ink/termio/dec.js'
\end{lstlisting}

Claude Code 对 Ink 渲染器进行了封装，添加了终端控制序列支持（如光标显示/隐藏）。

\section{JSX 与工具 UI}

工具可以通过 \code{setToolJSX} 设置自定义的 React UI：

\begin{lstlisting}[language=TypeScript,breaklines=true]
export type SetToolJSXFn = (
  args: {
    jsx: React.ReactNode | null
    shouldHidePromptInput: boolean
    shouldContinueAnimation?: true
    showSpinner?: boolean
    isLocalJSXCommand?: boolean
    isImmediate?: boolean
    clearLocalJSX?: boolean
  } | null,
) => void
\end{lstlisting}

这意味着每个工具执行时可以渲染自己的 UI 组件——例如，文件编辑工具可以显示差异视图，搜索工具可以显示实时结果。

\section{React Hooks 在 CLI 中的应用}

\code{src/hooks/} 目录包含大量 React Hooks：

\begin{lstlisting}[language=TypeScript,breaklines=true]
src/hooks/
├── toolPermission/      # 工具权限相关 Hooks
├── fileSuggestions.ts   # 文件建议
├── useCanUseTool.ts     # 工具使用权限
└── ...
\end{lstlisting}

例如，\code{useCanUseTool} Hook 在 REPL 组件中使用，每次工具调用都通过它进行权限检查。

\section{全屏界面 (Screens)}

\begin{lstlisting}[language=TypeScript,breaklines=true]
src/screens/
├── Doctor/      # /doctor 诊断界面
├── Resume/      # /resume 恢复界面
└── ...
\end{lstlisting}

某些命令（如 \code{/doctor}）需要全屏 UI。这些通过独立的 Screen 组件实现，接管整个终端。

\section{FPS 追踪}

\begin{lstlisting}[language=TypeScript,breaklines=true]
import type { FpsMetrics } from './utils/fpsTracker.js'

// 启动遥测中记录
lastFpsAverage: fpsMetrics?.averageFps,
lastFpsLow1Pct: fpsMetrics?.low1PctFps,
\end{lstlisting}

Claude Code 追踪 UI 的渲染帧率，记录平均 FPS 和 1% 最低 FPS。这反映了 Anthropic 对终端 UI 流畅度的重视。

\section{输出样式}

\begin{lstlisting}[language=TypeScript,breaklines=true]
src/outputStyles/   # 输出样式配置
\end{lstlisting}

支持不同的输出样式——包括主题、颜色方案等，可通过 \code{/theme} 命令切换。

\section{Vim 模式}

\begin{lstlisting}[language=TypeScript,breaklines=true]
src/vim/           # Vim 模式实现
\end{lstlisting}

Claude Code 甚至内置了 Vim 模式，通过 \code{/vim} 命令切换。这反映了其目标用户群——习惯在终端中工作的开发者。

\section{Buddy（伴侣精灵）}

\begin{lstlisting}[language=TypeScript,breaklines=true]
src/buddy/         # 伴侣精灵（彩蛋）
\end{lstlisting}

一个有趣的彩蛋——Claude Code 包含一个"伴侣精灵"组件，为终端体验增添趣味性。

\section{设计启示}

\subsection{React 的通用性}

React 不仅仅是一个"Web 框架"——它的核心是一个\textbf{声明式 UI 编程模型}。通过 \code{react-reconciler}，这个模型可以适配到任何渲染目标：浏览器 DOM、原生移动端、VR、甚至终端。

\subsection{组件化的终端 UI}

140+ 个组件的规模证明了组件化方法在复杂 CLI 应用中的可行性。传统的字符串拼接方式在这种规模下将变得不可维护。

\subsection{性能追踪}

FPS 追踪在 CLI 应用中很少见，但 Claude Code 的实时流式输出场景确实需要关注渲染性能。

\bigskip\noindent\rule{\textwidth}{0.4pt}\bigskip

\textbf{下一章}：\href{./chapter-13-feature-flags.md}{Feature Flag 与条件编译}
